\documentclass[11pt,a4paper]{ctexart}

\usepackage[margin=2.3cm]{geometry}
\usepackage{amsmath,amssymb,amsthm,mathtools}
\usepackage[dvipsnames]{xcolor}
\usepackage{hyperref}
\usepackage{enumitem}
\usepackage[most]{tcolorbox}
\usepackage{booktabs}

\hypersetup{colorlinks=true,linkcolor=blue!55!black,urlcolor=blue!55!black}
\setlist[enumerate]{itemsep=0.4em, topsep=0.5em}
\setlist[itemize]{itemsep=0.4em, topsep=0.5em}

\theoremstyle{definition}
\newtheorem{goalbox}{目标}
\newtheorem{attempt}{尝试}
\newtheorem{problem}{问题}
\newtheorem{definitionx}{定义}
\newtheorem{lemma}{引理}
\newtheorem{theorem}{命题}

\theoremstyle{remark}
\newtheorem*{remarkx}{备注}
\newtheorem*{intuitionx}{直觉}
\newtheorem*{summary}{小结}

\tcbset{
  commonbox/.style={enhanced, breakable,
    boxrule=0.8pt, arc=2mm,
    left=1.4mm, right=1.4mm, top=1mm, bottom=1mm, colback=white}
}
\newtcolorbox{keypoint}{commonbox,colframe=ForestGreen!60!black,colback=ForestGreen!4,
  title=本节交付,fonttitle=\bfseries}

\newcommand{\R}{\mathbb{R}}
\newcommand{\softmax}{\mathrm{softmax}}
\newcommand{\Att}{\mathrm{Att}}
\newcommand{\MH}{\mathrm{MultiHead}}

\title{Attention 的数学推导\\\large 像数学课本一样，从「想做什么」出发，一步一步推到 QKV 与多头}
\author{}
\date{\today}

\begin{document}
\maketitle

\begin{abstract}
本讲义不抢图、不堆术语，而是\emph{严格按照数学课本风格}走一遍 attention 的推导：
先把要解决的目标写清楚，写出建模假设，
然后从最朴素的尝试开始，
\textbf{每改进一步都说明「为什么必须这么改」}，
直到自然推出 Transformer 论文里那个完整的多头注意力公式。

读完之后，你应该能在白板上不查资料地写出
$\Att(Q,K,V) = \softmax(QK^\top / \sqrt{d_k})V$
的每一项是怎么来的、为什么这么来。
\end{abstract}

\tableofcontents
\newpage

%====================================================================
\section{第 0 步：先把目标写清楚}

我们手上有的对象很简单：

\begin{definitionx}[输入]
一段长度为 $n$ 的文本，每个 token 已经被嵌入成 $d$ 维向量。
把这些向量按行堆起来得到矩阵
\[
  X \;=\; \begin{pmatrix} x_1^\top \\ x_2^\top \\ \vdots \\ x_n^\top \end{pmatrix}
  \;\in\; \R^{n \times d}.
\]
$x_i \in \R^d$ 是第 $i$ 个 token 的输入向量。
\end{definitionx}

\begin{goalbox}[本章要做的事]
设计一个函数 $f: \R^{n \times d} \to \R^{n \times d}$，把输入 $X$ 变成输出
\[
  Y \;=\; f(X) \;\in\; \R^{n \times d},
\]
\textbf{满足}：
\begin{enumerate}
  \item \textbf{上下文感知}：$y_i$（输出的第 $i$ 行）必须能依赖\emph{所有} $x_1, \dots, x_n$，
        而不只是 $x_i$；
  \item \textbf{置换等变}（除位置编码外）：先打乱再算等价于先算再打乱；
  \item \textbf{可微}：$f$ 关于参数处处可导，能用反向传播训练；
  \item \textbf{并行}：所有 $y_i$ 必须能在 GPU 上一次性算出来，而不是循环。
\end{enumerate}
\end{goalbox}

\begin{remarkx}
\textbf{条件 1} 排除了「$y_i = g(x_i)$」这种逐位置网络（FFN）；
\textbf{条件 4} 排除了 RNN（必须按顺序算）。
满足这 4 条的最简单方案，就是 attention。
\end{remarkx}

%====================================================================
\section{第一次尝试：直接用相似度加权}

\subsection{动机}

最朴素的想法：要让 $y_i$ 看到所有 $x_j$，就把它们\emph{加起来}。
但「均匀加和」会丢失结构，所以要加权。
权重应该反映「$x_j$ 对 $x_i$ 有多重要」。

\textbf{怎么量化「重要」？} 在向量空间里最常用的相似度就是\emph{点积}：
两个向量越像，点积越大。

\begin{attempt}[无参数注意力]
对每对位置 $(i, j)$，定义未归一化分数
\[
  e_{ij} \;=\; x_i^\top x_j \;\in\; \R.
\]
对每一行做 softmax 把分数变成「合法权重」：
\[
  \alpha_{ij} \;=\; \frac{\exp(e_{ij})}{\sum_{k=1}^{n} \exp(e_{ik})}, \qquad \sum_{j=1}^{n} \alpha_{ij} = 1.
\]
最终输出是加权求和：
\[
  y_i \;=\; \sum_{j=1}^{n} \alpha_{ij}\, x_j.
\]
\end{attempt}

\subsection{矩阵化}

把上面写成矩阵形式：
\[
  E \;=\; X X^\top \;\in\; \R^{n \times n}, \qquad
  A \;=\; \softmax_{\text{行}}(E), \qquad
  Y \;=\; A X.
\]

逐项验证：
\begin{itemize}
  \item $E_{ij} = x_i^\top x_j$ ✓
  \item $A_{ij} = \alpha_{ij}$ ✓（softmax 沿每一行做）
  \item $(AX)_{i,:} = \sum_j A_{ij}\, x_j^\top = y_i^\top$ ✓
\end{itemize}

\subsection{它满足目标的几条吗？}

\begin{itemize}
  \item 上下文感知 ✓（每个 $y_i$ 都聚合了所有 $x_j$）
  \item 置换等变 ✓（公式里没有显式位置）
  \item 可微 ✓（softmax 与矩阵乘都是可导的）
  \item 并行 ✓（一次矩阵乘搞定）
\end{itemize}

\textbf{看起来全都满足了。那为什么还要 QKV？}\\
下一节回答这个问题。

%====================================================================
\section{为什么必须引入 $W_Q, W_K, W_V$}

\subsection{尝试 1 的两个根本缺陷}

\begin{problem}[对称性导致角色混淆]
注意尝试 1 里的相似度矩阵 $E = X X^\top$ 是\emph{对称的}：
\[
  E_{ij} = x_i^\top x_j = x_j^\top x_i = E_{ji}.
\]
但语言里「我看你」和「你看我」常常不该一样重要：
比如代词 \emph{it} 在英文里要指向某个名词，但那个名词不一定要回看 \emph{it}。
对称的相似度无法表达这种\emph{非对称的关注}。
\end{problem}

\begin{problem}[「相似度」与「贡献」被强行绑死]
在尝试 1 里，被加权的对象 $x_j$ 既参与\emph{打分}（出现在 $E_{ij} = x_i^\top x_j$），
又参与\emph{被聚合}（出现在 $\sum_j \alpha_{ij} x_j$）。
这意味着「$x_j$ 让自己显眼的方式」和「$x_j$ 真正贡献的内容」必须由\emph{同一个向量}承担。

\textbf{这是不必要的限制：}有时一个 token 的「让人找到我用什么标签」
（key）和「我能贡献什么内容」（value）应当是\emph{两件事}。
\end{problem}

\begin{problem}[模型没有任何可调旋钮]
$Y = \softmax(XX^\top) X$ 里\emph{没有任何参数}。
模型完全无法根据任务\emph{学到}「在这个任务里，什么算重要」。
\end{problem}

\subsection{解药：把每个 token 投影成三个不同的角色}

为了同时解决以上三个问题，最经济的改动是：
\textbf{给每个 $x_i$ 派三份不同的「替身」}，分别承担三种职能。

\begin{definitionx}[Query / Key / Value 投影]
引入三个可学习矩阵
\[
  W_Q \in \R^{d \times d_k},\quad W_K \in \R^{d \times d_k},\quad W_V \in \R^{d \times d_v}.
\]
对每个 token 定义
\[
  q_i \;=\; W_Q^\top x_i,\qquad k_j \;=\; W_K^\top x_j,\qquad v_j \;=\; W_V^\top x_j.
\]
矩阵化：$Q = X W_Q,\; K = X W_K,\; V = X W_V$，
形状 $Q, K \in \R^{n \times d_k}$，$V \in \R^{n \times d_v}$。
\end{definitionx}

\textbf{每个矩阵的物理意义}（这是理解 QKV 的关键）：
\begin{itemize}
  \item $W_Q$：「我以什么方式\emph{询问}别人」——把 $x_i$ 投到\emph{查询空间}；
  \item $W_K$：「我以什么方式\emph{被找到}」——把 $x_j$ 投到\emph{钥匙空间}；
  \item $W_V$：「如果真的有人来取，我贡献什么」——把 $x_j$ 投到\emph{内容空间}。
\end{itemize}

\subsection{重新定义打分函数}

把尝试 1 里的 $e_{ij} = x_i^\top x_j$ 换成
\begin{equation}
  e_{ij} \;=\; q_i^\top k_j \;=\; (W_Q^\top x_i)^\top (W_K^\top x_j) \;=\; x_i^\top W_Q W_K^\top x_j.
  \label{eq:score}
\end{equation}

\begin{lemma}[非对称性已经获得]
一般情况下 $W_Q W_K^\top \neq W_K W_Q^\top$，所以
\[
  e_{ij} = x_i^\top (W_Q W_K^\top) x_j \;\neq\; x_j^\top (W_Q W_K^\top) x_i = e_{ji}.
\]
即 \textbf{打分矩阵 $E = QK^\top$ 不再强制对称}，
模型可以学到「我看你 $\neq$ 你看我」的关注模式。
\end{lemma}

\begin{lemma}[「钥匙」和「内容」已经分离]
打分用的是 $k_j = W_K^\top x_j$，
聚合用的是 $v_j = W_V^\top x_j$。
两个投影矩阵\emph{独立}，所以「让别人找到我用什么标签」
和「我贡献什么」是\emph{两份独立的可学习内容}。
\end{lemma}

\subsection{得到 attention 的雏形}

把上面三件事写到一起：
\[
  E = Q K^\top \in \R^{n\times n}, \qquad
  A = \softmax_{\text{行}}(E), \qquad
  Y = A V \in \R^{n \times d_v}.
\]

合成一个公式：
\begin{equation}
  \boxed{\;\Att(Q, K, V) \;=\; \softmax(QK^\top)\, V.\;}
\end{equation}

到这一步还差一个细节：分母 $\sqrt{d_k}$。
下一节回答它。

%====================================================================
\section{为什么要除以 $\sqrt{d_k}$}

\subsection{问题：维度变大会让 softmax「过激」}

\begin{lemma}[点积的方差随维度增长]
设 $q, k \in \R^{d_k}$ 的各分量独立同分布，均值 0，方差 1。则
\[
  \mathbb{E}[q^\top k] \;=\; 0,\qquad \mathrm{Var}(q^\top k) \;=\; \sum_{i=1}^{d_k} \mathrm{Var}(q_i k_i) \;=\; d_k.
\]
\end{lemma}

也就是说，当 $d_k = 64$ 时，分数的标准差大约是 $8$；
当 $d_k = 512$ 时，标准差大约是 $\sqrt{512} \approx 22.6$。
分数越大，softmax 出来的概率分布越\emph{尖}（接近 one-hot）。

\subsection{尖锐分布的副作用：梯度消失}

设 softmax 的输入为 $z = (z_1, \dots, z_n)$，输出为 $p$，则
\[
  \frac{\partial p_i}{\partial z_j} \;=\; p_i (\delta_{ij} - p_j).
\]

当某个 $p_i \to 1$、其它 $\to 0$ 时，
$p_i(1 - p_i) \to 0$，所有偏导都趋于 0。
\textbf{反向传播时梯度几乎为零，模型学不动}。

\subsection{解药：把分数除以 $\sqrt{d_k}$}

把分数归一化到与 $d_k$ 无关的尺度：
\[
  \tilde{e}_{ij} \;=\; \frac{q_i^\top k_j}{\sqrt{d_k}}, \qquad
  \mathrm{Var}(\tilde e_{ij}) \;=\; \frac{d_k}{(\sqrt{d_k})^2} = 1.
\]

无论 $d_k$ 是 64 还是 512，分数的尺度都稳定在 $O(1)$，
softmax 不会过早饱和。

\subsection{完整的「缩放点积注意力」}

把缩放装回去：

\begin{theorem}[Scaled Dot-Product Attention]
\[
  \boxed{\;\Att(Q, K, V) \;=\; \softmax\!\left(\frac{QK^\top}{\sqrt{d_k}}\right) V.\;}
\]
\end{theorem}

\begin{keypoint}
到此为止我们已经把 §3.3 与 §3.4 全部推完了：
\textbf{每一步都不是天上掉下来的}——
QKV 是为了打破对称性、分离 key 与 value，并引入参数；
$\sqrt{d_k}$ 是为了 softmax 在不同维度下都能保持温度。
\end{keypoint}

%====================================================================
\section{为什么需要「多头」}

\subsection{单头的局限}

单头注意力的输出（不算 $W_V$）维度是 $n \times d_v$。
设想我们让 $d_v = d$（保持模型宽度不变）。
此时 $W_Q, W_K \in \R^{d \times d_k}$。

\textbf{关键观察：}注意力矩阵 $A = \softmax(QK^\top/\sqrt{d_k})$ 的形状是 $n \times n$，
也就是「每个位置看哪里」是\emph{一份}权重表。
\textbf{一份权重表只能编码一种关注模式。}

但语言里每个 token 同时存在多种关系：
\begin{itemize}
  \item 主语 $\to$ 谓语；
  \item 形容词 $\to$ 它修饰的名词；
  \item 代词 $\to$ 指代对象；
  \item 邻居关系 / 句法依存 / 同义对……
\end{itemize}

我们希望模型能\emph{同时}学若干份独立的关注模式。

\subsection{把「想要 \texorpdfstring{$h$}{h} 套独立关注」精确化}

「同时学 $h$ 种关注模式」翻译成数学，就是要 $h$ 张\emph{各不相同}的注意力矩阵：
\[
  A^{(1)}, A^{(2)}, \dots, A^{(h)} \;\in\; \R^{n \times n}.
\]
每张 $A^{(i)}$ 对应一种「看哪里」的规律。
对应地，我们也要 $h$ 段独立的输出
$\mathrm{out}^{(i)} = A^{(i)} V^{(i)}$，
最后还得有办法把这 $h$ 段拼成一段完整的输出。

\begin{remarkx}
\textbf{两个方案的共同点：}都给每个头一套独立的 $(W_Q^i, W_K^i, W_V^i)$，
都得到 $h$ 张独立的 $A^{(i)}$。\\
\textbf{两个方案的唯一区别：}\emph{每个头用多大的 $d_k$ 和 $d_v$}——也就是「每个头看的子空间有多宽」。
\end{remarkx}

\subsection{方案 A（笨办法）：每个头都开「全宽」}

最直接的做法：让\emph{每个头}都和原来的单头一样宽，即
\[
  d_k^{(A)} \;=\; d_v^{(A)} \;=\; d_{\text{model}}.
\]
形状：
\[
  W_Q^i,\, W_K^i,\, W_V^i \;\in\; \R^{d_{\text{model}} \times d_{\text{model}}}, \qquad i = 1, \dots, h.
\]
每个头：
\[
  \mathrm{head}_i^{(A)} \;=\; \Att(X W_Q^i,\, X W_K^i,\, X W_V^i)
  \;\in\; \R^{n \times d_{\text{model}}}.
\]

\textbf{代价核算}（投影矩阵的总参数量）：
\[
  \underbrace{3 d_{\text{model}}^2}_{\text{单头}} \;\to\;
  \underbrace{h \cdot 3 d_{\text{model}}^2}_{\text{$h$ 个全宽头}}.
\]
当 $h = 8, d_{\text{model}} = 512$ 时，光是 QKV 投影就从 $\approx 0.79\,$M 涨到 $\approx 6.3\,$M，\emph{每层都要乘以 8}。
计算量同样涨 $8$ 倍。\textbf{贵到不能接受。}

\begin{summary}
方案 A：「我想要 $h$ 个头」就字面上叠 $h$ 个。
做了同样的事 $h$ 次，参数和算量都变 $h$ 倍。
\end{summary}

\subsection{方案 B（聪明办法）：把同一份预算切成 \texorpdfstring{$h$}{h} 段}

\textbf{关键观察：}方案 A 之所以贵，是因为它\emph{不肯把预算切开}。
如果我们承认「单头能用的参数预算就是 $3 d_{\text{model}}^2$」，那么
最自然的折中是：

\begin{quote}
\emph{把这 $3 d_{\text{model}}^2$ 的预算\textbf{切成 $h$ 份}，每个头分到 $1/h$。
每个头看到的子空间维度 $d_k$ 也按比例缩小到 $d_{\text{model}}/h$。}
\end{quote}

\begin{definitionx}[多头注意力]
设 $d_k = d_v = d_{\text{model}} / h$。每个头的形状变成：
\[
  W_Q^i,\, W_K^i \;\in\; \R^{d_{\text{model}} \times d_k},\qquad
  W_V^i \;\in\; \R^{d_{\text{model}} \times d_v}.
\]
每个头：
\[
  \mathrm{head}_i \;=\; \Att\big(X W_Q^i,\; X W_K^i,\; X W_V^i\big)
  \;\in\; \R^{n \times d_v}, \qquad i = 1, \dots, h.
\]
然后拼接 + 一次输出投影：
\[
  \boxed{\;
  \MH(X) \;=\; \mathrm{Concat}(\mathrm{head}_1, \dots, \mathrm{head}_h)\, W_O,
  \quad W_O \in \R^{d_{\text{model}} \times d_{\text{model}}}.
  \;}
\]
拼接后的形状是 $n \times (h\cdot d_v) = n \times d_{\text{model}}$，
$W_O$ 把它映回 $d_{\text{model}}$ 维。
\end{definitionx}

\subsection{方案 A vs 方案 B 一图看完}

\begin{center}
\renewcommand{\arraystretch}{1.35}
\begin{tabular}{@{}p{4.2cm} c c@{}}
\toprule
                                  & \textbf{方案 A（全宽）}              & \textbf{方案 B（切片，Transformer 用的）} \\
\midrule
$d_k = d_v$                       & $d_{\text{model}}$                  & $d_{\text{model}}/h$ \\
单个 $W_Q^i$ 的形状               & $d_{\text{model}} \times d_{\text{model}}$ & $d_{\text{model}} \times (d_{\text{model}}/h)$ \\
$h$ 套 QKV 总参数量                & $h \cdot 3 d_{\text{model}}^2$       & $3 d_{\text{model}}^2$ \\
$h$ 张注意力矩阵 $A^{(i)}$         & 有 ✓                                & 有 ✓ \\
每张 $A^{(i)}$ 表达力              & 在\emph{完整} $d_{\text{model}}$ 维比对  & 只在 $d_k = d_{\text{model}}/h$ 维子空间比对 \\
$h$ 头是否独立学不同关注           & 是                                  & 是 \\
拼接后维度                         & $h \cdot d_{\text{model}}$           & $d_{\text{model}}$ \\
是否需要 $W_O$ 融合                & \emph{需要}（且必须把 $h\cdot d_{\text{model}}$ 压回去） & 需要（$d_{\text{model}}\to d_{\text{model}}$） \\
和单头比，参数额外开销             & $\approx 3(h-1) d_{\text{model}}^2 + h d_{\text{model}}^2$ & 只多了 $W_O$，约 $d_{\text{model}}^2$ \\
\bottomrule
\end{tabular}
\end{center}

\textbf{关键结论：}方案 B 用「\emph{把每个头看的子空间变窄}」这一个让步，
换来「总参数和总算量与单头几乎一样」。
\textbf{我们没有放弃「$h$ 张独立注意力」这一核心目标}——
A 和 B 都拿到了 $h$ 张独立的 $A^{(i)}$，
区别只是\emph{每张 $A^{(i)}$ 是在多大的子空间里算出来的}。

\subsection{为什么「砍小子空间」不致命？}

读到这里你应该会有一个自然怀疑：

\begin{quote}
\emph{把 $d_k$ 从 512 砍到 64，每个头看到的信息变少了 8 倍，
难道整体表达力不也跟着崩了 8 倍？}
\end{quote}

不会。三条理由，按重要性排：

\begin{enumerate}
  \item \textbf{我们要的本来就不是「每个头看见所有维度」，
        而是「每个头专注一种关系」。}
        一种语法关系（比如「主语—谓语」）通常只需要少数几个维度就能编码，
        不需要每个头都把整段 $d_{\text{model}}$ 重新审一遍。
        让头变窄，反而\emph{逼}它专精。
  \item \textbf{$h$ 份子空间合起来是\emph{完整}的 $d_{\text{model}}$。}
        因为 $h \cdot d_v = d_{\text{model}}$，
        $h$ 个头各看一段不重合的子空间，
        合起来正好覆盖整个 $d_{\text{model}}$ 维。
        没有信息总量上的损失，只是\emph{分工}方式不同。
  \item \textbf{$W_O$ 把分工的产出再融合回去。}
        每个头的输出在拼接之后只是\emph{并排放在一起}，
        $W_O$ 让模型可以学到「在最终输出里按什么权重混这 $h$ 个头」，
        恢复跨子空间的信息互动。
\end{enumerate}

\begin{summary}
方案 A 是「用蛮力的钱解决问题」；
方案 B 是「承认预算固定，用\emph{分工}解决问题」。
两者目标完全一致（拿到 $h$ 张独立的注意力），
只是 B 在每个头的「视野宽度」上做了让步，
换回了「便宜得起」。Transformer 选 B。
\end{summary}

\subsection{参数量核算（B 的精确账）}

\begin{lemma}[B 的总参数量与单头几乎相同]
\[
  \sum_{i=1}^{h} \dim(W_Q^i)
  = h \cdot d_{\text{model}} \cdot \frac{d_{\text{model}}}{h}
  = d_{\text{model}}^2.
\]
$W_K, W_V$ 同理。
所以「单头用整个 $d_{\text{model}}$」与「$h$ 头各用 $d_k = d_{\text{model}}/h$」
\emph{主体投影的总参数量完全相同}，都是 $3 d_{\text{model}}^2$。
唯一额外的开销是输出投影 $W_O$，多出 $d_{\text{model}}^2$。
\end{lemma}

\begin{intuitionx}
所以多头（方案 B）的几何含义是：\textbf{把一份 $d_{\text{model}}$ 维的「整空间注意力」
切成 $h$ 份「子空间注意力」，每份独立学一种关系，最后用 $W_O$ 把它们融合回完整空间。}\\
QKV 的\emph{机制}没有变化，只是\emph{在哪段子空间上做}从「整段」变成「$h$ 段」。
\end{intuitionx}

\subsection{为什么这等价于「reshape + transpose」}

工程实现里我们\emph{不会}真的开 $h$ 套小矩阵循环。
利用一个数学等价性：

\begin{lemma}[$h$ 个小投影 = 1 个大投影 + 切块]
把 $h$ 套 $W_Q^i \in \R^{d_{\text{model}} \times d_k}$ 沿\emph{列}方向拼起来，
得到一个大矩阵
\[
  W_Q \;=\; [W_Q^1, W_Q^2, \dots, W_Q^h] \;\in\; \R^{d_{\text{model}} \times (h\cdot d_k)} = \R^{d_{\text{model}} \times d_{\text{model}}}.
\]
则
\[
  X W_Q \;=\; [X W_Q^1,\; X W_Q^2,\; \dots,\; X W_Q^h]
  \;\in\; \R^{n \times d_{\text{model}}},
\]
也就是「一次大矩阵乘」的结果，
\textbf{逐列切开正好是 $h$ 段独立的 $\mathrm{head}_i$ 的 query}。
\end{lemma}

\textbf{这一引理才是「reshape + transpose」的真正源头}：
\begin{enumerate}
  \item 把 $h$ 套小投影写成一个大投影 $W_Q$，一次矩阵乘得到 $X W_Q \in \R^{n \times d_{\text{model}}}$；
  \item 把最后一维\emph{重塑}成 $(h, d_k)$，得到 $\R^{n \times h \times d_k}$；
  \item 转置让 $h$ 维移到前面，作为新的 batch 轴：$\R^{h \times n \times d_k}$。
\end{enumerate}

到这一步，$h$ 个头各自的 $\mathrm{head}_i$ 就可以\emph{用一次带 batch 的矩阵乘}同时算完。

\begin{summary}
「reshape + transpose」不是什么炫技，它就是
\emph{把数学上等价于「$h$ 个独立小矩阵」的存储布局，
摆成 GPU 喜欢的「带 batch 的张量」}。
意义和数学定义没有任何区别。
\end{summary}

%====================================================================
\section{把整套推导收成一个清单}

\begin{enumerate}
  \item \textbf{目标}：构造 $f: \R^{n \times d} \to \R^{n \times d}$，让每个输出位置都能依赖整段输入，且可微 / 可并行。
  \item \textbf{第一次尝试}：$Y = \softmax(XX^\top) X$。
        满足 4 条目标，但有三个缺陷：对称、key/value 绑死、无参数。
  \item \textbf{修补三缺陷}：引入 $Q = XW_Q,\, K = XW_K,\, V = XW_V$。
        得到 $\Att(Q,K,V) = \softmax(QK^\top) V$。
  \item \textbf{尺度修正}：维度变大会让点积方差变 $d_k$，
        softmax 饱和、梯度消失。
        除以 $\sqrt{d_k}$ 把方差稳定在 $1$，得到
        $\Att(Q,K,V) = \softmax(QK^\top/\sqrt{d_k}) V$。
  \item \textbf{多头扩展}：一份注意力只能学一种关注模式。
        把 $d_{\text{model}}$ 切成 $h$ 份子空间，每份独立做 attention，最后拼接 + $W_O$。
        总参数量与单头几乎相同，只多了一个 $W_O$。
  \item \textbf{工程等价}：$h$ 套小投影 = 1 套大投影 + reshape + transpose，
        所以实现里只看到「一个大 Linear + reshape + 一次 batched matmul」。
\end{enumerate}

\begin{keypoint}
\textbf{Attention 的整套公式都不是「定义」，而是「为了同时满足若干约束被推出来的结果」。}
你只要记住每一步的「为什么不能少」，就可以在白板上自己重新推一遍：

\[
  X \;\xrightarrow{\;W_Q, W_K, W_V\;} \; Q, K, V
  \;\xrightarrow{\;QK^\top/\sqrt{d_k}\;} \; \text{分数}
  \;\xrightarrow{\;\softmax\;} \; A
  \;\xrightarrow{\;A V\;} \; \text{输出}
  \;\xrightarrow{\;W_O\;} \; Y.
\]
\end{keypoint}

\end{document}
